\section{Controllability definitions}
\label{sec:definitions}

These definitions are shared across internal and boundary control -- the
thesis does not split them, and neither do we. Here $H$ and $U$ are the state
and control spaces of \eqref{eq:abstract-system}, as fixed in
Section~\ref{sec:wellposedness}.

\begin{definition}[Controllability notions for \eqref{eq:abstract-system}]
  \label{def:controllability}
  Let $T > 0$. System \eqref{eq:abstract-system} is said to be:
  \begin{enumerate}[label=(\roman*)]
    \item \textbf{exactly controllable} in time $T$ if, for all
      $(y_0,y_1)\in H\times H$, there exists $u \in L^2(0,T;U)$ such that the
      solution satisfies $y(T;y_0,u) = y_1$;
    \item \textbf{exactly controllable to trajectories} in time $T$ if, for
      all $(y_0,\hat y_0)\in H\times H$ and $\hat u \in L^2(0,T;U)$, there
      exists $u \in L^2(0,T;U)$ such that
      $y(T;y_0,u) = y(T;\hat y_0,\hat u)$;
    \item \textbf{null controllable} (or exactly controllable to zero) in
      time $T$ if, for all $y_0 \in H$, there exists $u \in L^2(0,T;U)$ such
      that $y(T;y_0,u)=0$;
    \item \textbf{approximately controllable} in time $T$ if, for all
      $(y_0,y_1)\in H\times H$ and every $\varepsilon>0$, there exists
      $u \in L^2(0,T;U)$ such that $\|y(T;y_0,u)-y_1\|_H < \varepsilon$.
  \end{enumerate}
\end{definition}

Equivalently, defining the reachable-set map
$\Phi: L^2(0,T;U)\to H$, $u \mapsto y(T;y_0,u)$, and its range
$R(T) = \{y(T;y_0,u): u\in L^2(0,T;U)\} = \Phi\big(L^2(0,T;U)\big) \subset H$,
the four notions above correspond to
$R(T)=H$, $R(T) \ni y(T;\hat y_0,\hat u)$ for every choice, $0 \in R(T)$, and
$R(T)$ dense in $H$, respectively. ($\Phi$ is distinct from the completion
space $F$ introduced in Remark~\ref{rem:hum-existence} below -- different
symbol, different role.)

\begin{remark}
  When there is no restriction on the norm of the initial condition, the
  properties in Definition~\ref{def:controllability} hold \emph{globally};
  under such a restriction, controllability holds only \emph{locally}.
\end{remark}

\subsection*{Steady-state controllability (constrained-problem aside)}

For the constrained/steady-state problem of Section~\ref{sec:constrained-aside}
the thesis restates controllability more concretely for the
boundary-controlled equation and adds:

\begin{definition}
  System \eqref{eq:heat-boundary} is (exactly) controllable to
  $y_1 \in L^2(\Omega)$ if, for every $y_0 \in L^2(\Omega)$, there exist
  $T_{y_0}>0$ and $u \in L^2(\Sigma)$ such that the solution satisfies
  $y(T_{y_0}) = y_1$.
\end{definition}

\begin{definition}[Steady state]
  $\bar y$ is a steady state of \eqref{eq:heat-boundary} if $-\Delta\bar y = 0$
  in $\Omega$, $\bar y = \bar u$ on $\partial\Omega$.
\end{definition}

\begin{definition}[Steady-state controllability]
  Let $y_0,y_1$ be steady states. System \eqref{eq:heat-boundary} is
  exactly-steady-state controllable in time $T>0$ if there exists
  $u \in L^2(\Sigma)$ such that the solution satisfies $y(\cdot,0)=y_0$ and
  $y(\cdot,T)=y_1$.
\end{definition}
